\documentclass[12pt,a4paper]{article}
\usepackage[utf8]{inputenc}
\usepackage[spanish,german]{babel}
\usepackage[T1]{fontenc}
\usepackage{geometry}
\geometry{a4paper, margin=2.5cm}
\usepackage{titlesec}
\usepackage{enumitem}
\usepackage{multicol}
\usepackage{csquotes}

\titleformat{\section}
{\Large\bfseries}
{}
{0em}
{}[\titlerule]

\titleformat{\subsection}
{\large\bfseries}
{}
{0em}
{}

\setlength{\parindent}{0pt}
\setlength{\parskip}{0.5em}

\begin{document}

\title{\textbf{DIALOG: Hablar del Tiempo}}
\author{Spanisch A1/A2}
\date{}
\maketitle

\textbf{Tema:} Hablar del tiempo \\
\textbf{Nivel:} A1 \\
\textbf{Complejidad:} Einfach \\
\textbf{Palabras:} 756

\vspace{1em}

\textit{Nota: Se recomienda revisar primero el vocabulario antes de leer el diálogo.}

\vspace{1em}

\section*{Diálogo}

\textbf{Laura:} ¡Hola, Pedro! ¿Cómo estás?

\textbf{Pedro:} ¡Hola, Laura! Muy bien, gracias. ¿Y tú?

\textbf{Laura:} Bien, gracias. Oye, ¿qué tiempo hace hoy? Hace mucho calor, ¿no?

\textbf{Pedro:} Sí, hace mucho calor. Creo que hay treinta grados. Es un día muy soleado.

\textbf{Laura:} Sí, el sol está muy fuerte. Por eso llevo gafas de sol y una gorra.

\textbf{Pedro:} Buena idea. Yo también llevo gorra. ¿Te gusta el calor?

\textbf{Laura:} No mucho. Prefiero el tiempo fresco. Me gusta más cuando hace veinte o veintidós grados.

\textbf{Pedro:} A mí también me gusta más el tiempo fresco. El calor es muy agotador.

\textbf{Laura:} Sí, es verdad. ¿Qué tiempo hace normalmente en tu ciudad?

\textbf{Pedro:} En mi ciudad hace mucho calor en verano, pero en invierno hace frío. A veces llueve mucho en invierno.

\textbf{Laura:} ¿Y en primavera y otoño?

\textbf{Pedro:} En primavera hace buen tiempo. Hace sol y hace calor, pero no demasiado. En otoño hace fresco y a veces llueve.

\textbf{Laura:} Me gusta mucho la primavera. Es mi estación favorita. ¿Cuál es tu estación favorita?

\textbf{Pedro:} Mi estación favorita es el otoño. Me gusta cuando hace fresco y los árboles cambian de color.

\textbf{Laura:} Eso es muy bonito. Oye, ¿has visto el pronóstico del tiempo para mañana?

\textbf{Pedro:} No, no lo he visto. ¿Qué dice?

\textbf{Laura:} Dice que mañana va a llover. Va a hacer mucho viento también.

\textbf{Pedro:} Ah, qué lástima. Yo quería ir a la playa mañana.

\textbf{Laura:} Sí, es una pena. Pero el domingo va a hacer buen tiempo otra vez. Va a hacer sol y va a hacer calor.

\textbf{Pedro:} Perfecto. Entonces puedo ir a la playa el domingo. ¿Te gusta la lluvia?

\textbf{Laura:} Sí, me gusta cuando llueve un poco, pero no cuando llueve mucho tiempo. ¿Y a ti?

\textbf{Pedro:} A mí me gusta la lluvia. Me gusta escuchar el sonido de la lluvia. Es muy relajante.

\textbf{Laura:} Sí, es verdad. Es relajante. ¿Qué haces cuando hace mal tiempo?

\textbf{Pedro:} Cuando hace mal tiempo, me quedo en casa. Veo películas o leo un libro. ¿Y tú?

\textbf{Laura:} Yo también me quedo en casa. A veces cocino o escucho música. Me gusta estar en casa cuando llueve.

\textbf{Pedro:} Sí, es muy cómodo. Oye, ¿has notado que el clima está cambiando?

\textbf{Laura:} Sí, lo he notado. Hace más calor que antes. Los veranos son más calurosos.

\textbf{Pedro:} Es verdad. Y llueve menos. A veces hay sequía.

\textbf{Laura:} Sí, es un problema. Pero hoy hace un día perfecto para estar afuera.

\textbf{Pedro:} Sí, tienes razón. Aunque hace mucho calor, es un día bonito.

\textbf{Laura:} Exacto. Bueno, tengo que irme. ¡Que tengas un buen día!

\textbf{Pedro:} Gracias, tú también. ¡Hasta luego!

\textbf{Laura:} ¡Hasta luego!

\newpage

\section*{Vocabulario}

\begin{multicols}{2}
\begin{itemize}[leftmargin=*]
    \item \textbf{tiempo} - Wetter, Zeit
    \item \textbf{¿Qué tiempo hace?} - Wie ist das Wetter?
    \item \textbf{hace calor} - es ist heiß
    \item \textbf{hace frío} - es ist kalt
    \item \textbf{hace sol} - es ist sonnig
    \item \textbf{hace viento} - es ist windig
    \item \textbf{grados} - Grad (Temperatur)
    \item \textbf{soleado} - sonnig
    \item \textbf{gafas de sol} - Sonnenbrille
    \item \textbf{gorra} - Mütze, Kappe
    \item \textbf{fresco} - kühl, frisch
    \item \textbf{agotador} - anstrengend, ermüdend
    \item \textbf{verano} - Sommer
    \item \textbf{invierno} - Winter
    \item \textbf{llover} - regnen
    \item \textbf{llueve} - es regnet
    \item \textbf{primavera} - Frühling
    \item \textbf{otoño} - Herbst
    \item \textbf{estación} - Jahreszeit
    \item \textbf{pronóstico del tiempo} - Wettervorhersage
    \item \textbf{viento} - Wind
    \item \textbf{playa} - Strand
    \item \textbf{lluvia} - Regen
    \item \textbf{relajante} - entspannend
    \item \textbf{mal tiempo} - schlechtes Wetter
    \item \textbf{clima} - Klima
    \item \textbf{caluroso} - heiß, schwül
    \item \textbf{sequía} - Dürre
    \item \textbf{afuera} - draußen
\end{itemize}
\end{multicols}

\newpage

\section*{Preguntas de Comprensión}

\begin{enumerate}[leftmargin=*]
    \item ¿Qué tiempo hace hoy según el diálogo?
    \item ¿Cuántos grados hace aproximadamente?
    \item ¿Qué lleva Laura para protegerse del sol?
    \item ¿Qué temperatura prefiere Laura?
    \item ¿Qué tiempo hace normalmente en la ciudad de Pedro en verano?
    \item ¿Cuál es la estación favorita de Laura?
    \item ¿Cuál es la estación favorita de Pedro?
    \item ¿Qué dice el pronóstico del tiempo para mañana?
    \item ¿Qué quería hacer Pedro mañana?
    \item ¿Qué tiempo va a hacer el domingo?
    \item ¿A Pedro le gusta la lluvia?
    \item ¿Qué hace Pedro cuando hace mal tiempo?
    \item ¿Qué hace Laura cuando llueve?
    \item ¿Qué han notado Laura y Pedro sobre el clima?
    \item ¿Por qué Pedro no puede ir a la playa mañana?
    \item ¿Qué le gusta a Pedro del otoño?
    \item ¿Qué tiempo hace en primavera según Pedro?
    \item ¿Qué tiempo hace en otoño según Pedro?
    \item ¿Qué problema mencionan sobre el clima?
    \item ¿Cómo se despiden Laura y Pedro?
\end{enumerate}

\newpage

\section*{Verbo Irregular: IR (gehen)}

\textit{Nota: \enquote{vosotros/as} se usa solo en España, no en español sudamericano. Se incluye aquí para completitud.}

\begin{multicols}{2}

\subsection*{Presente}
\begin{itemize}[leftmargin=*]
    \item yo - voy
    \item tú - vas
    \item él/ella/usted - va
    \item nosotros/as - vamos
    \item vosotros/as - vais
    \item ustedes - van
    \item ellos/ellas - van
\end{itemize}

\subsection*{Pretérito Indefinido}
\begin{itemize}[leftmargin=*]
    \item yo - fui
    \item tú - fuiste
    \item él/ella/usted - fue
    \item nosotros/as - fuimos
    \item vosotros/as - fuisteis
    \item ustedes - fueron
    \item ellos/ellas - fueron
\end{itemize}

\subsection*{Pretérito Imperfecto}
\begin{itemize}[leftmargin=*]
    \item yo - iba
    \item tú - ibas
    \item él/ella/usted - iba
    \item nosotros/as - íbamos
    \item vosotros/as - ibais
    \item ustedes - iban
    \item ellos/ellas - iban
\end{itemize}

\columnbreak

\subsection*{Futuro Simple}
\begin{itemize}[leftmargin=*]
    \item yo - iré
    \item tú - irás
    \item él/ella/usted - irá
    \item nosotros/as - iremos
    \item vosotros/as - iréis
    \item ustedes - irán
    \item ellos/ellas - irán
\end{itemize}

\subsection*{Pretérito Perfecto}
\begin{itemize}[leftmargin=*]
    \item yo - he ido
    \item tú - has ido
    \item él/ella/usted - ha ido
    \item nosotros/as - hemos ido
    \item vosotros/as - habéis ido
    \item ustedes - han ido
    \item ellos/ellas - han ido
\end{itemize}

\end{multicols}

\newpage

\section*{Explicación Gramatical}

\textbf{Das Verb \enquote{ir} -- Gehen und Zukunft}

Das Verb \enquote{ir} (gehen) ist eines der wichtigsten und häufigsten Verben im Spanischen. Es ist unregelmäßig und wird nicht nur für Bewegung verwendet, sondern auch für die Bildung des Futurs mit \enquote{ir + a + Infinitiv}.

\textbf{Verwendung von \enquote{ir}:}
\begin{itemize}[leftmargin=*]
    \item \textbf{Bewegung}: \enquote{Voy a la playa} -- Ich gehe zum Strand
    \item \textbf{Zukunft}: \enquote{Va a llover} -- Es wird regnen (wörtlich: Es geht regnen)
    \item \textbf{Pläne}: \enquote{Voy a estudiar} -- Ich werde studieren / Ich gehe studieren
\end{itemize}

\textbf{Die Struktur \enquote{ir + a + Infinitiv} für die Zukunft:}
Diese Konstruktion entspricht dem deutschen Futur (werden + Infinitiv) und wird sehr häufig verwendet, oft häufiger als das Futur Simple.

\textbf{Beispiele aus dem Text:}
\begin{itemize}[leftmargin=*]
    \item \enquote{Dice que mañana va a llover} -- Es sagt, dass es morgen regnen wird
    \item \enquote{Va a hacer mucho viento también} -- Es wird auch sehr windig sein
    \item \enquote{Yo quería ir a la playa mañana} -- Ich wollte morgen zum Strand gehen
    \item \enquote{Puedo ir a la playa el domingo} -- Ich kann am Sonntag zum Strand gehen
    \item \enquote{Va a hacer buen tiempo otra vez} -- Es wird wieder gutes Wetter sein
    \item \enquote{Va a hacer sol y va a hacer calor} -- Es wird sonnig und heiß sein
\end{itemize}

\textbf{Wichtige Unterschiede:}
\begin{itemize}[leftmargin=*]
    \item \enquote{ir + a + Infinitiv} = nahe Zukunft, Pläne, Absichten (z.B. \enquote{Voy a estudiar} -- Ich werde studieren)
    \item \enquote{Futur Simple} = allgemeine Zukunft, Vermutungen (z.B. \enquote{Estudiaré} -- Ich werde studieren)
    \item In der Umgangssprache wird \enquote{ir + a + Infinitiv} häufiger verwendet als das Futur Simple
\end{itemize}

\textbf{Weitere Verwendungen:}
\begin{itemize}[leftmargin=*]
    \item \enquote{ir + gerundio} = fortschreitende Handlung (z.B. \enquote{Voy mejorando} -- Ich werde besser)
    \item \enquote{ir + participio} = Zustand (z.B. \enquote{Voy vestido de azul} -- Ich bin blau gekleidet)
\end{itemize}

\end{document}
